\documentclass[a4paper,12pt,dutch]{article}
\usepackage{glossaries}
\usepackage[T1]{fontenc}
\usepackage{babel}
\usepackage{graphicx}
\usepackage[table,xcdraw]{xcolor}
\usepackage{hyperref}
\usepackage{blindtext}
\usepackage{geometry}
\usepackage{parskip}
\usepackage{mathtools}
\usepackage{siunitx}
\usepackage{listings}
\usepackage{csquotes}
\usepackage{caption}
\usepackage{subcaption}
\usepackage{comment}
\usepackage{pdfpages}
\usepackage{pict2e}

% \DeclareRobustCommand{\slashcirc}{{\mathpalette\doslashcirc\relax}}

% \makeatletter
% \newcommand\doslashcirc[2]{%
%   \sbox\z@{$#1\m@th\circ$}%
%   \setlength\unitlength{\wd\z@}
%   \begin{picture}(1,1)
%   \roundcap
%   \put(0,0){\box\z@}
%   \put(0,0){\line(1,1){1}}
%   \end{picture}%
% }
% \makeatother


% %% Some packages you will need
% \usepackage{pgfplots}
% \usepackage{pgfplotstable}
% \usepackage{booktabs}
% \usepackage{array}
% \usepackage{colortbl}


\definecolor{arduinoorange}{HTML}{FFA500}
\definecolor{arduinogray}{HTML}{808080}
\definecolor{arduinoblue}{HTML}{007ACC}
\definecolor{arduinogreen}{HTML}{469B00}

\lstset{
  language=C++,
  basicstyle=\ttfamily\footnotesize,
  keywordstyle=\color{arduinoorange},
  stringstyle=\color{arduinogreen},
  commentstyle=\color{arduinogray},
  moredelim=[s][\color{arduinoblue}]{\#}{\ },
  morekeywords={digitalRead,digitalWrite,pinMode,analogRead,analogWrite,Serial,begin,HIGH,LOW},
  frame=tb,
  tabsize=4,
  showstringspaces=false,
  breaklines=true,
  numbers=left,
  numberstyle=\tiny\color{arduinogray},
  numbersep=5pt,
  extendedchars=true,
  literate={á}{{\'a}}1 {ã}{{\~a}}1 {é}{{\'e}}1,
}

\lstdefinestyle{Arduino}
{
  language=C++,
  basicstyle=\ttfamily\footnotesize,
  keywordstyle=\color{arduinoorange},
  stringstyle=\color{arduinogreen},
  commentstyle=\color{arduinogray},
  moredelim=[s][\color{arduinoblue}]{\#}{\ },
  morekeywords={digitalRead,digitalWrite,pinMode,analogRead,analogWrite,Serial,begin},
  frame=tb,
  tabsize=4,
  showstringspaces=false,
  breaklines=true,
  numbers=left,
  numberstyle=\tiny\color{arduinogray},
  numbersep=5pt,
  extendedchars=true,
  literate={á}{{\'a}}1 {ã}{{\~a}}1 {é}{{\'e}}1,
  backgroundcolor=\color{black!85},
  rulecolor=\color{arduinoorange},
  frame=single,
  frameround=tttt,
  framexleftmargin=6pt,
  framexrightmargin=6pt,
  framextopmargin=6pt,
  framexbottommargin=6pt,
  breaklines=true,
  postbreak=\raisebox{0ex}[0ex][0ex]{\ensuremath{\color{red}\hookrightarrow\space}},
}

\usepackage[
    backend=biber,
    backref=true,
    backrefstyle=none,
    sortcites=true,
    sorting=none,
    doi=false, % doi informatie wordt niet weergegeven
    %uniquename=true,
    %uniquelist=true,
    maxcitenames=3,
    %issn=false, werkt niet
    language=american
]{biblatex}
\addbibresource{information/Bronnen.bib}
\DefineBibliographyStrings{dutch}{
    backrefpage = {blz.},
    backrefpages = {blz.},
}
\makeglossaries
\definecolor{Grey1}{HTML}{343434}
\graphicspath{{./Media/Figuren/}}
 \geometry{
 a4paper,
 total={170mm,257mm},
 left=20mm,
 top=20mm,
 }
\hypersetup{
    colorlinks=true,
    linkcolor=blue,
    filecolor=magenta,      
    urlcolor=cyan,
    pdftitle={Overleaf Example},
    pdfpagemode=FullScreen,
    }


\begin{document}
\title{
\includegraphics[width=3.5in]{img/Logo.png} \\
\vspace*{1in}
\textbf{Project 4}\\
\textit{Plan of Approach}\\
Version 1
}
\author{
\vspace*{0.5in} \\
  Written by:\\
  Laurens van der Drift\\
  Justin van der Reijden\\
  Luuk van Kappel\\
  Marnix Harmsen\\
		\vspace*{0.2in} \\
		The Hague University of Applied Sciences\\
        \textbf{Electrical Engineering}\\
        Delft, The Netherlands
       } 
\maketitle
\phantomsection
\section*{Versie Historie} \addcontentsline{toc}{section}{Versie Historie}

\begin{table}[h]
\begin{tabular}{|l|l|l|l|}
\hline
\rowcolor[HTML]{4472C4} 
{\color[HTML]{FFFFFF} \textbf{Versie}} &
  {\color[HTML]{FFFFFF} \textbf{Datum}} &
  {\color[HTML]{FFFFFF} \textbf{Wijzigingen}} &
  {\color[HTML]{FFFFFF} \textbf{Auteur}} \\ \hline
\rowcolor[HTML]{D9E1F2} 
1.0 &
  \multicolumn{1}{c|}{\cellcolor[HTML]{D9E1F2}05-02-2024} &
 N.v.t. &
  Alset Innovations \\ \hline

% \rowcolor[HTML]{FFFFFF} 
% 2.0 &
%   \multicolumn{1}{c|}{\cellcolor[HTML]{FFFFFF}7-4-2023} &
%  Feedback van FeedbackFruits toegepast &
%   Infra   Vroom \\ \hline

% \rowcolor[HTML]{D9E1F2} 
% 3.0 &
%   \multicolumn{1}{c|}{\cellcolor[HTML]{D9E1F2}4-6-2023} &
%  Bijgewerkt voor Assessment 3 &
%   Infra   Vroom \\ \hline

\end{tabular}
\end{table}
%\addcontentsline{toc}{section}{Verklarende Woordenlijst}
\printglossaries
\newglossaryentry{tender}
{
    name=\textit{tender},
    description={De inschrijving om op een kavel een wind turbine park te bouwen.}
}

\phantomsection
\addcontentsline{toc}{section}{Lijst van figuren}
\listoffigures

\newpage
\tableofcontents
% \include{chapters/11_Verdeling}
\phantomsection
\addcontentsline{toc}{section}{Referenties}
\printbibliography
\end{document}
